\section{De las pantallas exactas a la interfaz física de teoría M}
\label{mdp4:sec:interfaz-teoria-m-rev8}
\index{teoría M!interfaz física}
\index{dualidad T}
\index{teoría de cuerdas}
\index{CFT}

La primera parte construyó, en dominios matemáticos explícitos, las
reducciones
\[
12\longrightarrow11\longrightarrow10,
\qquad
26\longrightarrow24,
\]
y la involución de momento, enrollamiento y radio
\[
(m,w,R)\longleftrightarrow(w,m,R^{-1}).
\]
Estos resultados no son todavía una teoría física de cuerdas ni una acción
de supergravedad. Constituyen las pantallas y equivalencias que una
realización dimensional debe respetar.

La interfaz se tipa mediante un mapa declarado
\[
\mathcal R_M:\mathcal X_{\rm HMT}^{\rm enr}
\longrightarrow
\mathcal X_{11}^{\rm phys},
\]
cuyo codominio debe portar, como mínimo, el campo métrico \(g_{MN}\), la
tres-forma \(C_{MNP}\), el gravitino \(\psi_M\), su acción y las
transformaciones de supersimetría. Una especialización a cuerdas exige
además identificar radio físico, tensión, sectores de momento y
enrollamiento y el espacio de estados sobre el que se realiza la dualidad
T. Una lectura conforme o CFT requiere su propia álgebra de operadores,
gradación y correspondencia estado--operador.

La conexión positiva con la rama excepcional no es una flecha desnuda
Monstruo\(\to\)teoría M. Ambas prolongaciones conservan datos del mismo
estado enriquecido ---fase, gradación, orientación, costura y memoria---,
pero emplean funtores distintos. En la rama excepcional esos datos llegan a
\(V^\natural\), el Monstruo y Moonshine; en la rama dimensional alimentan
las pantallas y, sólo mediante \(\mathcal R_M\), una posible realización de
campos y cuerdas.

\begin{definition}[Admisibilidad interna de una compactificación]
\label{mdp4:def:admisibilidad-compactificacion-hmt-rev3}
Una compactificación candidata se representa por una terna
\[
 (x_\infty,\kappa,\mathcal U),
\]
donde \(x_\infty\in\mathcal X_{\rm HMT}^{\rm enr}\) es una sección
superviviente,
\[
 \kappa:\mathcal X_{\rm HMT}^{\rm enr}\dashrightarrow\mathcal M
\]
es el mapa, eventualmente parcial, hacia un espacio de módulos declarado
\(\mathcal M\), y
\(\mathcal U\) transporta métrica, orientación, registro y datos de frontera.
Definimos
\[
 \operatorname{Adm}_{\rm HMT}(x_\infty,\kappa,\mathcal U)=1
\]
si concurren las cuatro condiciones siguientes:
\begin{enumerate}
 \item la sección posee prolongación compatible a toda profundidad;
 \item las dualidades declaradas son compatibles en el punto considerado:
       \[
        \kappa(\mathsf T_{\rm HMT}x_\infty)
        =\mathsf T_{\mathcal M}(\kappa(x_\infty)),
       \]
       cuando ambos miembros están definidos;
 \item el espectro desciende al cociente de compactificación;
 \item los observables conservan sus dominios, codominios y tipos bajo
       \(\mathcal U\).
\end{enumerate}
\end{definition}

Este predicado es un criterio interno de admisibilidad y no una
identificación automática con el \emph{Swampland} físico. Los criterios de
distancia, torre, gravedad débil y ausencia de simetrías globales, así como
las condiciones geométricas de una compactificación concreta, se evalúan
después sobre \(\operatorname{im}\kappa\) y requieren los correspondientes
mapas métricos, campos y escalas. La definición anterior fija la interfaz que
debe preservar una realización; no demuestra por sí sola la existencia de
una compactificación física ni una conjetura de Swampland.

\begin{statusbox}[title={Alcance de la interfaz}]
Las pantallas, los cocientes y la involución T-dual son resultados
matemáticos exactos. El diccionario completo hacia
\((g_{MN},C_{MNP},\psi_M)\), la acción, la supersimetría y una CFT física es
una interfaz tipada que debe verificarse en cada realización. Esta frontera
no reduce el resultado matemático ni autoriza a presentarlo como dinámica
física ya cerrada.
\end{statusbox}
